\item
\emph{The subsets @@M0@@ of @@M1@@ correspond bijectively to the families @@M2@@ of classes of invertible sheaves on the fibers of @@M3@@.}
\item \emph{A family @@M4@@ of @@M5@@ is quasi-compact if and only if the corresponding family @@M6@@ is bounded.}
\end{enumerate}

Indeed, (i) follows from (3.1)\,; (ii), from the following fact: if @@M7@@ is a point of @@M8@@ and @@M9@@ is an algebraically closed extension of @@M10@@, then the @@M11@@-points of @@M12@@ correspond bijectively to the isomorphism classes of invertible sheaves on @@M13@@.

As for (iii), suppose that the subset @@M14@@ of @@M15@@ is quasi-compact. Then @@M16@@ is contained in an open subset @@M17@@ of finite type, and the inclusion @@M18@@ corresponds to an invertible sheaf @@M19@@ on an @@M20@@, where @@M21@@ is a surjective, (flat) morphism of finite type. Finally, @@M22@@ is of finite type, so @@M23@@ bounds @@M24@@.

Conversely, suppose that @@M25@@ is bounded by a coherent sheaf @@M26@@ on an @@M27@@, with @@M28@@ of finite type over @@M29@@. One may suppose @@M30@@ invertible (1.13 (i)) so that @@M31@@ defines a morphism @@M32@@. The image of @@M33@@ contains @@M34@@ and is quasi-compact, hence noetherian by (i). Consequently, @@M35@@ is quasi-compact, and we have (iii).

By 3.2, the usual definition of morphisms of finite type between two Picard schemes generalizes to the following:

\sgaStableTarget{sga6.e13.def.3-3}\textbf{Definition 3.3.} \emph{Let @@M36@@ and @@M37@@ be proper over noetherian @@M38@@. A morphism between the Picard functors}
@@M39@@
